\section{Introduction}
The present abundance of baryonic matter and known particle properties raise a retrospective question: can our rest masses, inventories, and calculations, together with current cosmological abundances, reconstruct how many baryons and antibaryons were present during the annihilation epoch, and how does that reconstruction compare with established knowledge? The cosmic microwave background and primordial nucleosynthesis constrain the small surviving net baryon excess \citep{Fields2025,Planck2020}. This excess and the total pair population are different quantities.

A hadron resonance gas (HRG) uses the rest masses, degeneracies, and charges of particle species to sum their populations and thermodynamic contributions \citep{Vovchenko2019}. These masses were already used in our calculations. Here, \emph{particle signatures} means these identifying catalog properties, rather than original spectral signals measured by this study. The model links such properties, a specified temperature, and equilibrium populations mathematically. Its usefulness at low temperature has been compared with lattice quantum chromodynamics, with agreement depending on the observable and proximity to the transition \citep{Bazavov2012}.

Annihilation removes equal numbers of baryons and antibaryons and preserves a pre-existing excess when baryon number is conserved. Inverse reactions also contribute to thermal evolution. \citet{Satarov2013} studied both contributions together with expansion and emphasized the importance of heavy antibaryon resonances above approximately 100 MeV. A useful comparison therefore separates equilibrium population reconstruction, the information supplied by current abundances, and the reaction history requiring kinetic calculation.

\textbf{The primary objective was to test whether particle-species properties and the present baryon abundance can reconstruct baryon and antibaryon numbers under annihilation-epoch conditions, and to compare that reconstruction with published results and approximations.} An inverse constraint linking the observed cosmic excess to the HRG baryon chemical potential was solved at eight temperatures. We also examined whether the present abundance alone uniquely determines an initial population. Complementary checks compared three public inventories, the literature's nucleon approximation, and a previous entropy reference \citep{Garcia2026}. The working hypothesis was that reconstruction requires an explicit temperature and thermal conditions in addition to masses and present abundance.
